\section{CONCLUSION}

The central hypothesis that a CNN can learn to extract altitude-correlated visual cues from monocular images of water surfaces, in a manner analogous to human visual distance estimation was validated by the experimental results. The WaterNet baseline achieved an MAE of 1.787~m and an $R^2$ of 0.551 on real-world data, despite having been trained entirely on synthetic imagery. Critically, all four CNN models produced lower RMSE than the onboard LiDAR sensor (2.337--3.591~m vs.\ 4.364~m), which suffered from frequent signal loss (20.9\% near-zero readings) and catastrophic outliers (maximum error of 77.1~m) due to the specular and semi-transparent properties of the water surface. But the ResNet50 based models have very low explicability score to be considered instead of the custom made WaterNet. This result confirms the practical motivation of the work: over water, a vision-based estimator provides predictable and bounded error behavior, and can suit as an input for flight control over water.

An unexpected finding was that the custom WaterNet architecture, trained from scratch on domain-specific data, outperformed the ResNet50-based models that leveraged ImageNet-pretrained features. This suggests that for tasks requiring discrimination based on subtle texture-scale and brightness-distribution differences within a single visual category, domain-specific feature learning with appropriately sized convolutional kernels ($5\times5$) is more effective than transferring deep representations optimized for semantic classification across thousands of object categories. The handcrafted feature vector proved most beneficial when paired with the ResNet50 backbone, where it reduced the error standard deviation by more than half and shifted the real-world bias from $-1.455$~m to $+0.118$~m, effectively compensating for the backbone's domain mismatch. The WaterNet-FusionLite variant, with only 310{,}K parameters, demonstrated that a competitive accuracy--complexity trade-off is achievable, making embedded deployment on constrained UAV platforms a realistic prospect.

The work also revealed persistent challenges that must be addressed before operational deployment. A systematic bias was observed across all models, originating from the brightness distribution mismatch between synthetic and real domains: synthetic images at low altitudes appear slightly brighter than their real-world counterparts, and this discrepancy propagates through the Value-channel preprocessing into the learned representations. Furthermore, prediction variance increases substantially at higher altitudes, where the visual differences between adjacent altitude levels become increasingly subtle as wave texture averages out. The ground-truth measurements themselves carry non-negligible uncertainty from the barometric-GPS fusion sensor, which limits the interpretability of the reported error metrics as true model accuracy.

\subsection{Future Work}

Several directions are identified for extending this research toward a deployable system. First, the sim-to-real domain gap should be addressed through brightness histogram matching between synthetic and real training samples, and through domain adaptation techniques or further domain randomization of rendering parameters (lighting angle, atmospheric haze, water turbidity). Second, the systematic bias observed in all models motivates the development of a post-processing calibration stage, such as spline-based or piecewise-linear correction fitted on a small set of real-world calibration samples. Third, the real-world dataset should be expanded with measurements from diverse water bodies (rivers, lakes, reservoirs) under varying weather and illumination conditions, using a calibrated radar altimeter as ground truth rather than the barometric-GPS fusion sensor employed in this study. Fourth, model explainability should be strengthened through Grad-CAM analysis on the CNN models and deeper analysis on the feature fusion branches, to verify that the learned attention patterns correspond to physically meaningful regions of the water surface. Fifth, the deployment pipeline should be implemented as a ROS~2 node with TensorRT or ONNX quantization for real-time inference on embedded platforms such as the NVIDIA Jetson Orin Nano, including latency profiling and integration with the flight controller's altitude fusion filter. Also, a multi-output branch can help with a classification task, inferring if the subject is readable as water or not, to flag the controller to avoid using the measures outside the water surface and possibly confusing the sensor fusion filter.